% 第二章 系统总体设计

\section{系统总体设计}

\subsection{系统需求分析}
\par 本系统面向家庭环境监测与便捷交互场景，目标是在可控成本与可维护实现的前提下，构建一套覆盖“端侧采集—家庭网关—平台服务—交互应用”的完整闭环。需求分析以“可用性优先、工程可验证”为原则，将系统目标拆分为功能需求、性能需求与非功能需求三类。
\par 功能需求方面，系统需要对温度、湿度、气压、VOC及气体相关指标等环境信息进行周期采集，并在本地与平台侧提供可视化展示。系统需要具备端侧人脸采集、识别、名单管理与图像获取能力，以支撑家庭访问控制或状态感知等扩展应用。系统需要提供网页端交互入口，实现仪表板展示、历史数据查询、参数配置与人脸系统的远程控制。系统还需要提供语音交互入口，使用户可通过自然语言完成环境信息查询等任务，并通过工具调用机制把意图映射到受控接口以保障可控性。
\par 性能需求方面，系统的数据链路需要满足家庭场景的“秒级刷新、可感知实时性”。环境采集与显示应以独立节拍运行，避免某一模块阻塞导致全局停顿。控制链路需要具备可感知的响应速度，并在网络波动或设备重启时能够恢复工作。数据准确性以传感器规格与合理的采样策略为基础，同时通过异常值过滤、缺失值处理等方式提升数据可解释性。
\par 非功能需求方面，系统应具备良好的可扩展性，便于后续增加传感器类型、拓展控制对象或替换平台接口。系统应具备基本安全能力，至少在平台侧实现用户登录鉴权与接口权限控制，避免未授权访问控制接口。系统应具备可维护性与可移植性，支持在常见开发环境下复现构建与运行，并通过模块化分层减少耦合度。

\subsection{系统架构设计}
\par 为满足上述需求，本文采用端边云协同的分层架构。整体上可划分为硬件层、通信层与应用层三部分。硬件层由STM32采集节点、ESP32-S3家庭网关、OpenMV视觉节点与本地显示构成，用于完成环境感知、边缘视觉与就地显示。通信层负责端到端数据与命令的可靠传输，涵盖STM32与网关之间的UART链路、网关到平台的HTTP数据上报链路，以及平台到网关的MQTT命令与响应链路。应用层由Node.js服务器与Web前端组成，提供数据接收、鉴权、存储与可视化，同时为语音助手提供统一的数据接口。
\par 模块划分遵循“职责单一、接口稳定”的原则。采集侧模块完成传感器初始化、采样与数据封装；网关侧模块完成串口接收解析、MQ气体传感器采样、状态与数据结构管理、TFT显示刷新、WiFi连接维护、HTTP上报与MQTT消息收发；视觉侧模块完成采集、识别与命令解析；平台侧模块完成数据管理、配置管理、用户管理与MQTT管理，并向前端提供REST风格接口。
\par 数据流向方面，环境数据由采集侧产生并汇聚至网关，在网关完成融合后以JSON形式通过HTTP上传至服务器，服务器写入本地数据文件并向前端提供查询接口；控制流由前端或平台服务发起，经服务器校验后通过MQTT向网关下发命令，网关再通过串口转发给OpenMV执行，OpenMV的响应（文本或图像）回传至网关并经MQTT返回服务器，最终在网页端展示。

\par 配图建议如下。图2-1建议为“系统分层架构图”，画法为三层横向布局：硬件层（STM32/ESP32-S3/OpenMV/TFT）、通信层（UART/HTTP/MQTT/WiFi）、应用层（Node.js服务器/Web前端/语音助手），并用双向箭头标注层间交互。图2-2建议为“数据流向与控制流示意图”，画法为两条并行流程：数据流（采集→网关→HTTP→服务器→前端）与控制流（前端/服务器→MQTT→网关→UART→OpenMV→回传），并用不同颜色或线型区分。图2-3建议为“模块划分与接口关系图”，画法为模块框图，展示采集侧/网关侧/视觉侧/平台侧/交互侧五大模块及其内部子模块，并用连线标注模块间接口。
\begin{figure}[H]
\centering
\includegraphics[width=0.92\textwidth]{paper/fig/2-1.png}
\caption{系统分层架构图}
\label{fig:system_layers}
\end{figure}
\begin{figure}[H]
\centering
\includegraphics[width=0.92\textwidth]{paper/fig/2-2.png}
\caption{数据流向与控制流示意图}
\label{fig:data_flow}
\end{figure}
\begin{figure}[H]
\centering
\includegraphics[width=0.92\textwidth]{paper/fig/2-3.png}
\caption{模块划分与接口关系图}
\label{fig:module_structure}
\end{figure}

\subsection{技术选型}
\par 技术选型以“可获得性、生态成熟度与教学可复现”为主要依据。硬件平台方面，采集侧选用STM32F103系列作为典型MCU平台，具备完善的外设资源与成熟的开发生态，便于驱动I\(^2\)C类传感器并进行稳定的数据上报。网关侧选用ESP32-S3，具备WiFi联网能力与较好的算力/存储资源，可同时承担协议栈、显示与消息转发任务，并可在Arduino生态与PlatformIO环境下快速实现工程化开发。视觉侧选用OpenMV作为轻量视觉平台，能够在资源受限条件下完成基于Haar级联的人脸检测与基于LBP特征的轻量识别，适合作为端侧视觉的原型验证平台。
\par 通信协议方面，设备间短距离通信采用UART，具有实现简单、调试直观与工程可控的特点，适合采集侧到网关的数据上报与网关到OpenMV的命令交互。网关到服务器的数据上报采用WiFi + HTTP方式，以JSON作为数据载体便于平台解析与扩展。跨端控制与异步消息采用MQTT发布/订阅模型，能够较好适应家庭网络波动，并便于将“命令”与“响应/事件”解耦管理。
\par 软件平台方面，服务器端选用Node.js + Express实现REST接口与静态页面托管，配合MQTT客户端实现设备控制链路；嵌入式端采用Arduino框架与相关库（WiFi、HTTP、MQTT、TFT显示等）降低工程实现复杂度；语音交互子系统采用Python实现，组合ASR、LLM与TTS模块，并以工具调用方式接入平台API，便于扩展更多“可执行”的智能体能力。
\par 开发工具方面，嵌入式工程统一采用PlatformIO + VS Code进行依赖管理、编译与烧录，服务器端采用pnpm进行依赖管理，语音助手采用Python虚拟环境与requirements进行依赖管理，从而保证项目可复现与版本可控。
